\IfFileExists{stacks-project.cls}{%
\documentclass{stacks-project}
}{%
\documentclass{amsart}
}

% For dealing with references we use the comment environment
\usepackage{verbatim}
\newenvironment{reference}{\comment}{\endcomment}
%\newenvironment{reference}{}{}
\newenvironment{slogan}{\comment}{\endcomment}
\newenvironment{history}{\comment}{\endcomment}

% For commutative diagrams we use Xy-pic
\usepackage[all]{xy}

% We use 2cell for 2-commutative diagrams.
\xyoption{2cell}
\UseAllTwocells

% We use multicol for the list of chapters between chapters
\usepackage{multicol}

% This is generally recommended for better output
\usepackage{lmodern}
\usepackage[T1]{fontenc}

% For cross-file-references
\usepackage{xr-hyper}

% Package for hypertext links:
\usepackage{hyperref}

% For any local file, say "hello.tex" you want to link to please
% use \externaldocument[hello-]{hello}
\externaldocument[introduction-]{introduction}
\externaldocument[conventions-]{conventions}
\externaldocument[sets-]{sets}
\externaldocument[categories-]{categories}
\externaldocument[topology-]{topology}
\externaldocument[sheaves-]{sheaves}
\externaldocument[sites-]{sites}
\externaldocument[stacks-]{stacks}
\externaldocument[fields-]{fields}
\externaldocument[algebra-]{algebra}
\externaldocument[brauer-]{brauer}
\externaldocument[homology-]{homology}
\externaldocument[derived-]{derived}
\externaldocument[simplicial-]{simplicial}
\externaldocument[more-algebra-]{more-algebra}
\externaldocument[smoothing-]{smoothing}
\externaldocument[modules-]{modules}
\externaldocument[sites-modules-]{sites-modules}
\externaldocument[injectives-]{injectives}
\externaldocument[cohomology-]{cohomology}
\externaldocument[sites-cohomology-]{sites-cohomology}
\externaldocument[dga-]{dga}
\externaldocument[dpa-]{dpa}
\externaldocument[sdga-]{sdga}
\externaldocument[hypercovering-]{hypercovering}
\externaldocument[schemes-]{schemes}
\externaldocument[constructions-]{constructions}
\externaldocument[properties-]{properties}
\externaldocument[morphisms-]{morphisms}
\externaldocument[coherent-]{coherent}
\externaldocument[divisors-]{divisors}
\externaldocument[limits-]{limits}
\externaldocument[varieties-]{varieties}
\externaldocument[topologies-]{topologies}
\externaldocument[descent-]{descent}
\externaldocument[perfect-]{perfect}
\externaldocument[more-morphisms-]{more-morphisms}
\externaldocument[flat-]{flat}
\externaldocument[groupoids-]{groupoids}
\externaldocument[more-groupoids-]{more-groupoids}
\externaldocument[etale-]{etale}
\externaldocument[chow-]{chow}
\externaldocument[intersection-]{intersection}
\externaldocument[pic-]{pic}
\externaldocument[weil-]{weil}
\externaldocument[adequate-]{adequate}
\externaldocument[dualizing-]{dualizing}
\externaldocument[duality-]{duality}
\externaldocument[discriminant-]{discriminant}
\externaldocument[derham-]{derham}
\externaldocument[local-cohomology-]{local-cohomology}
\externaldocument[algebraization-]{algebraization}
\externaldocument[curves-]{curves}
\externaldocument[resolve-]{resolve}
\externaldocument[models-]{models}
\externaldocument[functors-]{functors}
\externaldocument[equiv-]{equiv}
\externaldocument[pione-]{pione}
\externaldocument[etale-cohomology-]{etale-cohomology}
\externaldocument[proetale-]{proetale}
\externaldocument[relative-cycles-]{relative-cycles}
\externaldocument[more-etale-]{more-etale}
\externaldocument[trace-]{trace}
\externaldocument[crystalline-]{crystalline}
\externaldocument[spaces-]{spaces}
\externaldocument[spaces-properties-]{spaces-properties}
\externaldocument[spaces-morphisms-]{spaces-morphisms}
\externaldocument[decent-spaces-]{decent-spaces}
\externaldocument[spaces-cohomology-]{spaces-cohomology}
\externaldocument[spaces-limits-]{spaces-limits}
\externaldocument[spaces-divisors-]{spaces-divisors}
\externaldocument[spaces-over-fields-]{spaces-over-fields}
\externaldocument[spaces-topologies-]{spaces-topologies}
\externaldocument[spaces-descent-]{spaces-descent}
\externaldocument[spaces-perfect-]{spaces-perfect}
\externaldocument[spaces-more-morphisms-]{spaces-more-morphisms}
\externaldocument[spaces-flat-]{spaces-flat}
\externaldocument[spaces-groupoids-]{spaces-groupoids}
\externaldocument[spaces-more-groupoids-]{spaces-more-groupoids}
\externaldocument[bootstrap-]{bootstrap}
\externaldocument[spaces-pushouts-]{spaces-pushouts}
\externaldocument[spaces-chow-]{spaces-chow}
\externaldocument[groupoids-quotients-]{groupoids-quotients}
\externaldocument[spaces-more-cohomology-]{spaces-more-cohomology}
\externaldocument[spaces-simplicial-]{spaces-simplicial}
\externaldocument[spaces-duality-]{spaces-duality}
\externaldocument[formal-spaces-]{formal-spaces}
\externaldocument[restricted-]{restricted}
\externaldocument[spaces-resolve-]{spaces-resolve}
\externaldocument[formal-defos-]{formal-defos}
\externaldocument[defos-]{defos}
\externaldocument[cotangent-]{cotangent}
\externaldocument[examples-defos-]{examples-defos}
\externaldocument[algebraic-]{algebraic}
\externaldocument[examples-stacks-]{examples-stacks}
\externaldocument[stacks-sheaves-]{stacks-sheaves}
\externaldocument[criteria-]{criteria}
\externaldocument[artin-]{artin}
\externaldocument[quot-]{quot}
\externaldocument[stacks-properties-]{stacks-properties}
\externaldocument[stacks-morphisms-]{stacks-morphisms}
\externaldocument[stacks-limits-]{stacks-limits}
\externaldocument[stacks-cohomology-]{stacks-cohomology}
\externaldocument[stacks-perfect-]{stacks-perfect}
\externaldocument[stacks-introduction-]{stacks-introduction}
\externaldocument[stacks-more-morphisms-]{stacks-more-morphisms}
\externaldocument[stacks-geometry-]{stacks-geometry}
\externaldocument[moduli-]{moduli}
\externaldocument[moduli-curves-]{moduli-curves}
\externaldocument[examples-]{examples}
\externaldocument[exercises-]{exercises}
\externaldocument[guide-]{guide}
\externaldocument[desirables-]{desirables}
\externaldocument[coding-]{coding}
\externaldocument[obsolete-]{obsolete}
\externaldocument[fdl-]{fdl}
\externaldocument[index-]{index}

% Theorem environments.
%
\theoremstyle{plain}
\newtheorem{theorem}[subsection]{Theorem}
\newtheorem{proposition}[subsection]{Proposition}
\newtheorem{lemma}[subsection]{Lemma}

\theoremstyle{definition}
\newtheorem{definition}[subsection]{Definition}
\newtheorem{example}[subsection]{Example}
\newtheorem{exercise}[subsection]{Exercise}
\newtheorem{situation}[subsection]{Situation}

\theoremstyle{remark}
\newtheorem{remark}[subsection]{Remark}
\newtheorem{remarks}[subsection]{Remarks}

\numberwithin{equation}{subsection}

% Macros
%
\def\lim{\mathop{\mathrm{lim}}\nolimits}
\def\colim{\mathop{\mathrm{colim}}\nolimits}
\def\Spec{\mathop{\mathrm{Spec}}}
\def\Hom{\mathop{\mathrm{Hom}}\nolimits}
\def\Ext{\mathop{\mathrm{Ext}}\nolimits}
\def\SheafHom{\mathop{\mathcal{H}\!\mathit{om}}\nolimits}
\def\SheafExt{\mathop{\mathcal{E}\!\mathit{xt}}\nolimits}
\def\Sch{\mathit{Sch}}
\def\Mor{\mathop{\mathrm{Mor}}\nolimits}
\def\Ob{\mathop{\mathrm{Ob}}\nolimits}
\def\Sh{\mathop{\mathit{Sh}}\nolimits}
\def\NL{\mathop{N\!L}\nolimits}
\def\CH{\mathop{\mathrm{CH}}\nolimits}
\def\proetale{{pro\text{-}\acute{e}tale}}
\def\etale{{\acute{e}tale}}
\def\QCoh{\mathit{QCoh}}
\def\Ker{\mathop{\mathrm{Ker}}}
\def\Im{\mathop{\mathrm{Im}}}
\def\Coker{\mathop{\mathrm{Coker}}}
\def\Coim{\mathop{\mathrm{Coim}}}

% Boxtimes
%
\DeclareMathSymbol{\boxtimes}{\mathbin}{AMSa}{"02}

%
% Macros for moduli stacks/spaces
%
\def\QCohstack{\mathcal{QC}\!\mathit{oh}}
\def\Cohstack{\mathcal{C}\!\mathit{oh}}
\def\Spacesstack{\mathcal{S}\!\mathit{paces}}
\def\Quotfunctor{\mathrm{Quot}}
\def\Hilbfunctor{\mathrm{Hilb}}
\def\Curvesstack{\mathcal{C}\!\mathit{urves}}
\def\Polarizedstack{\mathcal{P}\!\mathit{olarized}}
\def\Complexesstack{\mathcal{C}\!\mathit{omplexes}}
% \Pic is the operator that assigns to X its picard group, usage \Pic(X)
% \Picardstack_{X/B} denotes the Picard stack of X over B
% \Picardfunctor_{X/B} denotes the Picard functor of X over B
\def\Pic{\mathop{\mathrm{Pic}}\nolimits}
\def\Picardstack{\mathcal{P}\!\mathit{ic}}
\def\Picardfunctor{\mathrm{Pic}}
\def\Deformationcategory{\mathcal{D}\!\mathit{ef}}
\title{The cosk\_0 fibre product of a surjection is a trivial Kan fibration}
\author{The Stacks Project Authors}
\date{stacks-project:simplicial.tex:local-snapshot-20261012; GNU FDL 1.2 or later}
\begin{document}
\maketitle
\noindent\textit{Editorial note.} The counted statement is the author's own lemma environment labelled lemma-cosk-minus-one-equivalence (source lines 6367-6378); the proof environment begins on the second line after it and is its own proof, closing at source line 6391.\par\medskip
\noindent\textit{Editorial note.} The statement and the proof invoke three author results of the same chapter by label; all three are retained as context spans and the generator expands the context transitively (budget 800 lines for this card), while reference targets in other chapters that carry canonical project tags are auto-declared as non-load-bearing source pointers.\par\medskip
\noindent\textit{Editorial note.} The retained preamble is the chapter's own minimal preamble (sources/stacks-simplicial/preamble.tex, taken with its real line count), a verbatim excerpt of the native Xy-pic setup, the native theorem declarations with the mathematical macros, and the project's own no-op reference environments; no macro was invented or redefined.\par\medskip
\noindent\textit{Editorial scope.} One original result of the Stacks Project's Simplicial Methods chapter, counted once: the cosk\_0 fibre product of a surjection is a trivial Kan fibration.\par\medskip
\noindent\textit{Allowed standard prerequisites (editorial).}\begin{itemize}
\item Simplicial sets and cosk\_0. The category of simplicial sets, the constant simplicial set on a set, and the coskeleton cosk\_0.. Both the constructed simplicial set and the target of the map are described through these (source lines 6369-6377, 6383-6386).
\item The fibre-product simplicial set. The simplicial set associated to a map of sets whose n-simplices are the (n+1)-fold fibre products (the author's example-fibre-products-simplicial-object, retained as context).. It is the simplicial set whose triviality is asserted (source line 6374).
\item Trivial Kan fibrations and base change. A surjection of sets gives a trivial Kan fibration cosk\_0(A) to cosk\_0(B) (the author's lemma-section), and trivial Kan fibrations are stable under base change (the author's lemma-trivial-kan-base-change), both retained as context.. They are the two steps of the proof (source lines 6388-6390).
\end{itemize}\medskip
\begin{example}
\label{example-fibre-products-simplicial-object}
Suppose that $Y\to X$ is a morphism of $\mathcal{C}$ such that all
the fibred products $Y \times_X Y \times_X \ldots \times_X Y$ exist.
Then we set $U_n$ equal to the $(n + 1)$-fold fibre product,
and we let $\varphi : [n] \to [m]$ correspond to the map
(on ``coordinates'')
$(y_0, \ldots, y_m) \mapsto (y_{\varphi(0)}, \ldots, y_{\varphi(n)})$.
In other words, the map $U_0 = Y \to U_1 = Y \times_X Y$ is the
diagonal map. The two maps $U_1 = Y \times_X Y \to U_0 = Y$ are the
projection maps.
\end{example}

\begin{lemma}
\label{lemma-section}
Let $f : V \to U$ be a morphism of simplicial sets. Let $n \geq 0$ be an
integer. Assume
\begin{enumerate}
\item The map $f_i : V_i \to U_i$ is a bijection for $i < n$.
\item The map $f_n : V_n \to U_n$ is a surjection.
\item The canonical morphism $U \to \text{cosk}_n \text{sk}_n U$
is an isomorphism.
\item The canonical morphism $V \to \text{cosk}_n \text{sk}_n V$
is an isomorphism.
\end{enumerate}
Then $f$ is a trivial Kan fibration.
\end{lemma}

\begin{proof}
Consider a solid diagram
$$
\xymatrix{
\partial \Delta[k] \ar[r] \ar[d] & V \ar[d] \\
\Delta[k] \ar[r] \ar@{-->}[ru] & U
}
$$
as in Definition \ref{definition-trivial-kan}. Let $x \in U_k$ be
the $k$-simplex corresponding to the lower horizontal arrow.
If $k \leq n$ then the dotted arrow is the one corresponding to
a lift $y \in V_k$ of $x$; the diagram will commute as the other
nondegenerate simplices of $\Delta[k]$ are in degrees $< k$ where
$f$ is an isomorphism. If $k > n$, then by conditions (3) and (4)
we have (using adjointness of skeleton and coskeleton functors)
$$
\Mor(\Delta[k], U) =
\Mor(\text{sk}_n\Delta[k], \text{sk}_nU) =
\Mor(\text{sk}_n\partial\Delta[k], \text{sk}_nU) =
\Mor(\partial \Delta[k], U)
$$
and similarly for $V$ because
$\text{sk}_n\Delta[k] = \text{sk}_n\partial\Delta[k]$ for $k > n$.
Thus we obtain a unique dotted arrow fitting into the diagram in
this case also.
\end{proof}

\begin{lemma}
\label{lemma-trivial-kan-base-change}
Let $f : X \to Y$ be a trivial Kan fibration of simplicial sets.
Let $Y' \to Y$ be a morphism of simplicial sets.
Then $X \times_Y Y' \to Y'$ is a trivial Kan fibration.
\end{lemma}

\begin{proof}
This follows immediately from the functorial properties of the fibre
product (Lemma \ref{lemma-fibre-product}) and the definitions.
\end{proof}

\begin{definition}
\label{definition-trivial-kan}
A map $X \to Y$ of simplicial sets is called a {\it trivial Kan fibration}
if $X_0 \to Y_0$ is surjective and for all $n \geq 1$ and any commutative
solid diagram
$$
\xymatrix{
\partial \Delta[n] \ar[r] \ar[d] & X \ar[d] \\
\Delta[n] \ar[r] \ar@{-->}[ru] & Y
}
$$
a dotted arrow exists making the diagram commute.
\end{definition}

\begin{lemma}
\label{lemma-fibre-product}
If $U, V, W$ are simplicial objects in the category $\mathcal{C}$,
and if $a : V \to U$, $b : W \to U$ are morphisms
and if $V \times_U W$ exists, then we have
$$
\Mor(T, V \times_U W) =
\Mor(T, V) \times_{\Mor(T, U)}
\Mor(T, W)
$$
for any fourth simplicial object $T$ of $\mathcal{C}$.
\end{lemma}

\begin{proof}
Omitted.
\end{proof}

\begin{lemma}
\label{lemma-cosk-minus-one-equivalence}
Let $A$, $B$ be sets, and that $f : A \to B$ is a map. Consider the simplicial
set $U$ with $n$-simplices
$$
A \times_B A \times_B \ldots \times_B A\ (n + 1 \text{ factors)}.
$$
see Example \ref{example-fibre-products-simplicial-object}.
If $f$ is surjective, the morphism $U \to B$
where $B$ indicates the constant simplicial set with value $B$
is a trivial Kan fibration.
\end{lemma}

\begin{proof}
Observe that $U$ fits into a cartesian square
$$
\xymatrix{
U \ar[d] \ar[r] & \text{cosk}_0(A) \ar[d] \\
B \ar[r] & \text{cosk}_0(B)
}
$$
Since the right vertical arrow is a trivial Kan fibration by
Lemma \ref{lemma-section}, so is the left by
Lemma \ref{lemma-trivial-kan-base-change}.
\end{proof}

\end{document}
